\section{税粮、钱币与家户：国家进入北方社会的路径}

\subsection{收入不是一个数字}

金朝的财政在史书中常以收入、盐课、田税、钱币和军费出现，仿佛这些项目可以汇成一个可以比较的总数。其实每一个数额都带着产生它的机关、时间、地域、币制和政治目的。中都的报告不等于一条边路的实收；规定征收不等于农户已经缴纳；某年仓粮数字也不等于当地居民实际拥有的粮食。若把不同文献的数字抽离口径拼成一条上升或下降曲线，得到的只是后人的整齐，不是金人面对的财政。

较值得追问的是财物如何移动。田赋必须经由里正、州县、仓场和道路进入更大的系统；盐、酒、矿冶或市场税需要检查、许可、运输和稽核；军费又会把粮、布、钱和劳力调到边防或都城。每一次转运都会产生新的裁量：地方官要决定催征的次序，仓吏要登记进出，商人要评估道路是否安全，农家要在交纳、储粮和借贷之间安排季节。国家的收入因此不是从土地自动流向宫廷，而是一条由人、文书与交通不断维持的链。

\subsection{南北边界的账目}

绍兴和议形成后，淮河与秦岭一带不只是军事线。两侧驻军、堡寨、渡口和使节往来都要消耗粮食和钱财，边市和私下贸易也会在禁令、许可与地方利益之间起伏。对金廷而言，南方边界需要从华北和中都得到补给；对南宋而言，守江淮亦需依靠东南财赋。所谓和平把成本从战场移到常态化的驻防、修城和转运，并未使边民脱离国家的计算。

一一六四年隆兴和议，宋金隆兴和议，恢复了外交框架，却不能保证每一处边地都按同一速度恢复交易或农业。河流的涨落、军队调防、盗匪、灾害和地方官的执行方式，都会让一个条约在不同渡口呈现不同效果。宋方诗文中关于“遗民”的痛苦，是理解南方政治记忆的重要入口，却不能替金境居民逐一说话；金方志书中的赋役条目也不能证明边地家庭没有逃避、协商或另谋生路。边界的两侧都存在看得见的命令与看不见的调整。

\subsection{钱币、实物与信用}

货币问题尤其不宜用现代“货币政策”的词汇轻易覆盖。铜钱、银、布帛、粮食和各类实物在不同地区、不同交易中可以并存；钱币的铸造、携带、成色和流通范围，会影响一笔税、一段工资或一次远途交易的实际价值。出土钱币能说明某类钱曾到达某地，窖藏的封藏期也许提示战争或不安，却不能直接告诉我们每个人怎样使用钱，更不能把一批遗物等同于全国的景气或崩溃。

市场中的信用关系同样难以从官书完整看见。商人可能凭熟人、家族、寺院或行会联系借贷和担保；农户在荒年借粮，来年以收成、劳力或田产偿付；军队经过时，地方可能被迫垫付，也可能从供应中获利。契约若能保存，会让某一次交涉突然具体，但保存地区和当事人本身就有偏差。没有契约不表示没有债务，有一份契约也不表示制度在全境相同。读者需要把宏观的财政名目与这种细小而有风险的安排放在一起。

\subsection{盐、铁与工匠的城市}

中都及华北城市的生产不能只按“官营”或“民营”二分。国家关心军器、盐、矿冶和都城供给，因而会设法控制关键环节；工匠、商人和运输者同时有自己的家庭、师徒、市场和迁徙路径。窑址、冶炼遗存和城市遗迹能显示生产曾集中在哪些空间，不能自动列出劳动者的工资、身份或情愿。若将物质遗存与制度文本、墓志或地方记载并读，才能看见国家需求如何进入作坊，以及工匠怎样在征调、雇佣和流动之间寻找余地。

中都的建设把这种关系放大。城墙、宫殿和官署需要材料，也需要把材料从不同区域调到城内。它创造了订单、运输和技术聚集的机会，同时可能把附近家户拖入劳役和物价波动。一个工程给都城留下的并非只有壮观遗迹，还包括道路磨损、木材来源、河运压力和地方村落的时间。正史较容易记录工程完成和大臣奏议；谁承担了哪一段搬运、谁因劳役延误耕作，通常只在偶然材料中闪现。

\subsection{土地并不只是税基}

金廷需要土地记录，因为税粮和军费需要稳定的预期；家庭需要土地，因为居住、婚姻、继承和抵御荒年都依赖它。两种需要重叠而不相同。登记可能把一块田变成可征收的面积，家户却同时关心灌溉、邻界、债务、坟墓和亲属的使用权。战争与迁徙会使旧边界失效，新的安置又会重绘关系。将田土单纯看成国家可计算的资源，便会忽略它为何能成为诉讼和亲属冲突的中心。

关于金代乡村的直接材料远少于我们希望得到的程度。后出的地方志会把较晚的地理和制度投回金代；墓志保存的多是有资源立石的家族；契约的地域分布也不平均。因此可以说，征赋、战争和人口移动改变了乡村条件，不能凭少数文本断言某种租佃或家族结构覆盖全朝。材料的沉默不应诱使作者虚构“普通农民”的统一声音。更诚实的写法，是让田地、道路、粮仓和家户各自留下可核的一部分，而将不能确认的普遍性留在问题之中。

\subsection{妇女、婚姻与继承的证据边缘}

婚姻、嫁资、寡居、收养和继承最能显示制度与家庭如何相互进入，却也最容易被后世道德语言遮蔽。法律或礼书所规定的秩序，不能直接等同于每个家庭的安排；墓志的妇德叙述也服务于纪念者的愿望。只有在文本日期、当事人关系和地方情境可辨时，我们才能谨慎地看到一次财产转移、一次家户重组或一位女性与寺院、亲属之间的联系。

战争年代使这一问题更尖锐。丈夫、兄弟或成年子弟被征发、俘虏或迁徙后，留在原地的人要处理田产、债务与儿童。材料往往用“户”“口”记录他们，而不是记录他们如何决定生计。不能因史书不写便说女性没有行动，也不能把一个墓志或传奇当成所有人的经验。正因为看不全，叙事应停止在证据所及之处：国家的登记如何把家庭列为单位，家庭又为何未必按这张表生活。

\subsection{寺院不是财政之外的岛屿}

寺院的功德碑、造像题记、经卷和土地关系，使宗教进入金代社会史。寺院可以是礼拜和修行的地方，也可能保存捐施、借贷、葬祭、教育和地方交往的痕迹。国家对僧道的管理、赏赐或限制，说明朝廷希望怎样分类宗教人员；它不能自动说明一座寺院的经济状况，更不能说明所有信众的内心。功德碑写的是捐施者希望被记住的善行，因而比起抽象的“民间信仰”，更适合让我们看见某次营建、某组姓名和某个地点。

多种文字材料尤其值得谨慎。女真文、契丹文与汉文题记的并存，表明政治和文化生活中存在不同书写传统；它不允许我们由一种文字推断稳定的民族身份或忠诚。翻译、转写和碑刻保存状态也会影响今天能读到什么。把这些局部材料放回中都、东北或边地的具体网络，才能避免把文化写成脱离税粮、交通和迁徙的装饰。

\begin{causalchain}{财政能力如何同时造成恢复与脆弱}
金朝接收华北后，需要以户籍、赋税、仓储、道路和军政编制维持都城与边防；这些安排使战争后的资源能够被重新调动。可是调动越依赖稳定的地方生产，战时征发、工程劳役和道路破坏就越可能削弱它。当蒙古战争迫使金廷多线应对时，原先连接中都、河南与边区的财政链既提供抵抗资源，也暴露出任何一段断裂都会传到全局的风险。
\end{causalchain}
